\begin{table*}[t]
\centering
\scriptsize
\setlength{\tabcolsep}{3pt}
\caption{R363 paper-facing visualization portfolio. The views are generated from tracked R320/R345/R348/R354/R355/R358 artifacts and summarize E2/E3 tradeoffs; R363 is not a new empirical result.}
\label{tab:visualization-portfolio}
\begin{tabular}{p{0.16\linewidth}p{0.19\linewidth}p{0.45\linewidth}p{0.09\linewidth}}
\toprule
View & Artifact & Paper-facing evidence & Role \\
\midrule
Baseline tradeoff & \texttt{baseline-tradeoff.svg} & Op-stack query Work@5 0.0937 vs flat 1; R@30\% 0.39 vs fixed-session 0.3559; groups 157.5 vs 285. & E2 \\
Metric heatmap & \texttt{metric-heatmap.svg} & Shows AP 0.3117, R@30\% 0.39, Work@5 0.0937, and the fixed-session WTFP counterpoint 0.0044. & E2 \\
Diagnostic lenses & \texttt{diagnostic-lenses.svg} & Six lenses over 36 objective rows: operation-stack-family views win 11/36, while baseline counterpoints win 25/36. & E3 \\
Actionability knobs & \texttt{actionability-knobs.svg} & 27/36 objective rows require non-default actions; R354 accepts 5/6 profile-spec patches; R358 boundary fields add AP 0.0181. & E3 \\
Oracle-depth adequacy & \texttt{oracle-depth-adequacy.svg} & Lower top-5 unit work than flat on 24/24 rows; higher fixed-session unit recall on 20/24; fewer groups to 50\% positives on 22/24. & E2/E3 \\
\bottomrule
\end{tabular}
\end{table*}
